% Pollis New Methods: nonlinear dynamics (threshold autoregression) — draft 1 (13 October 2026)
% Compile with: latexmk -pdf nm_threshold.tex && latexmk -c
% Shared look for the Pollis New Methods documents (same fonts and colours as the brochure and the workflow handouts).
\documentclass[10pt,a4paper]{article}

\usepackage[a4paper,left=18mm,right=18mm,top=15mm,bottom=17mm,footskip=8mm]{geometry}
\usepackage[T1]{fontenc}
\usepackage[default,semibold]{sourcesanspro}
\usepackage{textcomp}
\usepackage[italic]{mathastext}
\usepackage{amsmath}
\usepackage{microtype}
\usepackage[table]{xcolor}
\usepackage{tikz}
\usepackage[most]{tcolorbox}
\usepackage{fontawesome5}
\usepackage{tabularx,booktabs}
\usepackage{enumitem}
\usepackage{fancyhdr}
\usepackage[hidelinks]{hyperref}

% ---------------------------------------------------------------- colours (as in the brochure)
\definecolor{navy}{HTML}{1C3A5E}
\definecolor{pblue}{HTML}{3B78B5}
\definecolor{amber}{HTML}{F2B33D}
\definecolor{deepamber}{HTML}{D9902A}
\definecolor{pale}{HTML}{EEF3F9}
\definecolor{ink}{HTML}{2A2F36}
\definecolor{logonavy}{HTML}{173D63}
\definecolor{logodark}{HTML}{C98A1E}
\definecolor{logoamber}{HTML}{F2B446}
\definecolor{logocream}{HTML}{FEEFD2}
\definecolor{soft}{HTML}{5E6773}
\definecolor{agree}{HTML}{2F7D4F}
\definecolor{differ}{HTML}{B23B3B}
\color{ink}

\setlength{\parindent}{0pt}
\setlength{\parskip}{4pt}
\setlist{nosep,leftmargin=*,itemsep=2pt}
\renewcommand{\arraystretch}{1.15}

\pagestyle{fancy}
\fancyhf{}
\renewcommand{\headrulewidth}{0pt}
\fancyfoot[L]{\footnotesize\color{soft}\textbf{Pollis}\enspace New Methods\enspace·\enspace\nmtheme}
\fancyfoot[R]{\footnotesize\color{soft}\thepage}
\newcommand{\nmtheme}{}
\newcommand{\nmdate}{11 October 2026}

% The Pollis logo, redrawn so it prints sharply at any size
\newcommand{\pollislogo}[1]{%
  \begin{tikzpicture}[x=#1/512,y=-#1/512,baseline=(current bounding box.center)]
    \fill[logonavy] (0,0) rectangle (512,512);
    \fill[logodark] (36,422) -- (104,256) -- (172,422) -- cycle;
    \fill[logoamber] (104,422) -- (203,180) -- (302,422) -- cycle;
    \fill[logocream] (203,422) -- (339,89) -- (476,422) -- cycle;
  \end{tikzpicture}}

% Top of the first page: logo, name; then the kicker, title and a one-line promise
\newcommand{\nmhead}[3]{%
  \renewcommand{\nmtheme}{#1}%
  \thispagestyle{fancy}%
  \begin{minipage}[c]{13mm}\pollislogo{11mm}\end{minipage}%
  \begin{minipage}[c]{\dimexpr\linewidth-13mm\relax}
    {\color{navy}\Large\bfseries Pollis}\enspace{\color{soft}New Methods}\hfill
    {\color{deepamber}\small\bfseries\MakeUppercase{#1}}
  \end{minipage}\par\vspace{7mm}
  {\color{navy}\fontsize{22}{26}\selectfont\bfseries #2}\par\vspace{4pt}
  {\color{soft}\large #3}\par\vspace{5mm}}

\newcommand{\nmsection}[1]{\par\vspace{5pt}{\color{navy}\large\bfseries #1}\par\vspace{1pt}}
\newcommand{\nmsub}[1]{\par\vspace{2pt}{\color{deepamber}\bfseries #1}\par}
\newcommand{\gt}{\,{\color{deepamber}\bfseries ›}\,}
\newcommand{\menupath}[1]{{\bfseries\color{navy}#1}}
\newcommand{\fn}[1]{\texttt{\small #1}}
% Citations: \nmcite{key}{Breiman (2001)} links to the entry \nmref{key}{...} in the References
\newcommand{\nmcite}[2]{\hyperlink{ref:#1}{\color{pblue}#2}}
\newcommand{\nmref}[2]{\raisebox{1.1em}[0pt][0pt]{\hypertarget{ref:#1}{}}#2\par}
\newcommand{\nmdoi}[1]{\href{https://doi.org/#1}{\color{pblue}doi:#1}}
% Where a notebook is in Pollis: its title in Browse Methods, and its file
\newcommand{\browse}[2]{\menupath{Compose \gt New Methods \gt Browse Methods}: \emph{#1} (\fn{#2.ipynb})}

% Key facts at the top: source, missing in Julia, notebook, result
\newcommand{\factsheet}[4]{%
  \begin{tcolorbox}[enhanced, colback=pale, colframe=pale, boxrule=0pt, sharp corners,
    left=7pt, right=7pt, top=5pt, bottom=5pt, fontupper=\small]
  \begin{tabularx}{\linewidth}{@{}>{\leavevmode\bfseries\color{navy}}p{27mm}>{\leavevmode}X@{}}
  Tested against & #1\\
  Missing in Julia & #2\\
  Notebook & #3\\
  Result & #4
  \end{tabularx}
  \end{tcolorbox}}

% Boxes
\newtcolorbox{pbox}[2][]{enhanced, colback=pale, colframe=amber, boxrule=0pt,
  leftrule=3pt, sharp corners, left=7pt, right=7pt, top=6pt, bottom=6pt,
  fontupper=\small\raggedright, before upper={{\color{navy}\normalsize\bfseries #2}\par\smallskip}, #1}
\newtcolorbox{nbox}[2][]{enhanced, colback=navy, colframe=navy, boxrule=0pt,
  sharp corners, left=8pt, right=8pt, top=7pt, bottom=7pt, coltext=white,
  fontupper=\small\raggedright, before upper={{\color{amber}\normalsize\bfseries #2}\par\smallskip}, #1}

% Value tables: published beside Julia
\newenvironment{valuetable}[1]{%
  \par\small\tabularx{\linewidth}{#1}\toprule}{\bottomrule\endtabularx\par}
\newcommand{\ok}{{\color{agree}\faCheck}}
\newcommand{\no}{{\color{differ}\faTimes}}

% End of the document: credit and licence
\newcommand{\nmcredit}{%
  \par\vfill
  \begin{minipage}[b]{0.78\linewidth}\footnotesize\color{soft}
  Methods written and tested with Claude Code (Anthropic) in Pollis, from the published papers and their data.\par
  \textcopyright\ 2026 Antonio Saragga Seabra\enspace·\enspace Licensed under \href{https://creativecommons.org/licenses/by-sa/4.0/}{CC BY-SA 4.0}\enspace·\enspace Draft, \nmdate
  \end{minipage}\hfill
  \begin{minipage}[b]{0.18\linewidth}\raggedleft\pollislogo{9mm}\par\vspace{2pt}{\footnotesize contact@pollis.org}\end{minipage}}


\begin{document}

\nmhead{Nonlinear dynamics}
  {Threshold autoregression: SETAR models and limit cycles}
  {Why do the lynx rise for six years and crash in three? Tong and Lim (1980), reproduced.}

\factsheet
  {\nmcite{tong1980}{Tong and Lim (1980)}, \emph{Journal of the Royal Statistical Society: Series B}, eqs.\ 9.1--9.3 and
   Figs.\ 8, 10 and 12 (pp.\,259--262): the Canadian lynx trappings of 1821--1934 (PollisDatasets \fn{lynxmackenzie}) and
   Wolf's yearly sunspot numbers of 1700--1955 (R's \fn{sunspot.year}, read into memory and not saved).}
  {the self-exciting threshold autoregressive model (SETAR), its identification by AIC, and its eventual forecasting
   function and limit cycles. No registered Julia package fits it.}
  {\browse{Threshold autoregression: SETAR models and limit cycles}{ThresholdAR}}
  {68 of 70 printed values to every digit, the other two in the last. Every coefficient of the three models to 4 decimals.}

\nmsection{The problem}
Many cycles in nature and in economies are not symmetric: the Canadian lynx population climbs for about six years and
collapses in three; sunspot activity rises faster than it falls. A linear autoregression cannot produce this. Its cycles
are symmetric, and without noise it dies out to its mean. \nmcite{tong1980}{Tong and Lim (1980)} let the autoregression
change with the level of the series itself: one linear model below a threshold, another above it. So simple a
nonlinearity gives asymmetric cycles, limit cycles and jumps.

\nmsection{The method}
In SETAR$(2; k_1, k_2)$ with delay $d$ and threshold $r$, the series follows
\[ x_t = a_0^{(j)} + a_1^{(j)} x_{t-1} + \dots + a_{k_j}^{(j)} x_{t-k_j} + \varepsilon_t^{(j)}, \qquad
   j = 1 \text{ if } x_{t-d} \le r, \; j = 2 \text{ otherwise}, \]
and each regime is fitted by least squares to its own observations (\fn{setar}). \fn{identify} follows the paper's
Section 8: the candidate thresholds are the 30--70 per cent points of $x_{t-d}$, each regime's order is the one with the
smallest AIC, $n_j \ln(\mathrm{RSS}_j/n_j) + 2(k_j + 1)$, and the threshold the one with the smallest total.
\fn{selectdelay} compares delays by the normalised AIC, the AIC per observation fitted. \fn{eventual} continues the
series by the model without noise, the paper's eventual forecasting function; \fn{limitcycle} finds the cycle it settles
into and \fn{risedescent} the rise and descent of each local cycle. \fn{predict} gives one-step-ahead predictions.

\nmsection{What the tests show}
\begin{valuetable}{@{}lXl@{}}
\textbf{Test} & \textbf{What is compared} & \textbf{Agreement}\\\midrule
Eq.\,9.1 & lynx SETAR(2; 8, 3): 13 coefficients & every one to 4 decimals\\
Eq.\,9.2 & sunspots 1700--1920, SETAR(2; 4, 12): 18 coefficients & every one to 4 decimals\\
Eq.\,9.3 & sunspots 1837--1924, SETAR(2; 4, 2): 8 coefficients & every one to 4 decimals\\
p.\,259 & lynx limit cycle, period 9: 9 values & 8 to every digit, 1 in the last\\
p.\,261 & sunspot cycle: period 31, rise and descent of 3 local cycles & all 7\\
Eqs.\,9.1--9.3 & variances and NAICs: 12 values & 11 to every digit, 1 in the last\\
Figs.\,8--12 & lynx root mean RSS, two sunspot prediction RMSEs & 2 to every digit, 1 in the last\\
\end{valuetable}

The last-digit differences, side by side (Julia against the paper):
\begin{itemize}
  \item The lynx limit cycle: 3.22805 vs 3.2281.
  \item The NAIC of eq.\,9.3: 5.08178 vs 5.0817, which looks truncated rather than rounded.
  \item The one-step-ahead RMSE of eq.\,9.3 over 1925--1944: 9.23414 vs 9.2342.
\end{itemize}

\nmsection{Two things the paper does not say}
\nmsub{The variances and NAICs are those of a smaller model}
The residual variances printed under each equation, and the NAICs of the figures, are not those of the models as stated.
They are those of the same models with one lag fewer in each regime, and the NAIC penalises them as if they had the
orders stated. All three models show it, to every printed digit: eq.\,9.2's variances 254.64, 66.80 and 153.71 are those of
SETAR(2; 3, 11), and SETAR(2; 4, 12) gives 252.64, 66.55 and 152.65. The coefficients, the root mean RSS and the prediction
errors are those of the models as stated: the 1921--1955 RMSE is 12.17397 for SETAR(2; 4, 12), as printed, and 12.35812
for SETAR(2; 3, 11). It looks like an off-by-one in the program that printed the summaries.

\nmsub{The lynx cycle is printed backwards}
Read forwards, the printed lynx cycle rises for three years and falls for six; the text, and the model, rise for six
and fall for three. The paper lists the cycle from its smallest value backwards.

\nmsection{Still open}
The identification does not arrive at the published models by itself: for the lynx the smallest NAIC is at $d = 3$,
SETAR(2; 5, 9), and at $d = 2$ the search prefers SETAR(2; 10, 2). Tong and Lim chose with the eventual forecasting
function, the prediction errors of held-out observations and judgement as well (pp.\,258--259); \fn{selectdelay} returns
every candidate beside the one with the smallest NAIC. The published threshold itself, 3.1163, is on the grid: it is the
60 per cent point of $x_{t-2}$, the 1902 count of 1,307 lynx.

\nmsection{Further considerations}
\begin{itemize}
  \item \textbf{More regimes and other thresholds.} The paper's TARSO and TARSC models, with an outside series or several
    series in the threshold, need data the paper does not give in full; they would need a published table with its data
    before they become a new method here.
  \item \textbf{Inference on the threshold.} Testing linearity against SETAR, and confidence sets for $r$, came later
    (Chan, Tsay, Hansen); each would need its own published table to be tested against.
  \item \textbf{Upstream.} SETAR would sit naturally in a Julia time-series package; the three equations are a
    ready-made test for such a contribution.
\end{itemize}

\nmsection{References}
{\footnotesize
\nmref{tong1980}{Tong, H. and Lim, K.\,S. (1980). Threshold autoregression, limit cycles and cyclical data. \emph{Journal of the Royal Statistical Society: Series B} 42(3), 245--268. \nmdoi{10.1111/j.2517-6161.1980.tb01126.x}}
\nmref{elton1942}{Elton, C. and Nicholson, M. (1942). The ten-year cycle in numbers of the lynx in Canada. \emph{Journal of Animal Ecology} 11(2), 215--244.}
\nmref{campbell1977}{Campbell, M.\,J. and Walker, A.\,M. (1977). A survey of statistical work on the Mackenzie River series of annual Canadian lynx trappings for the years 1821--1934 and a new analysis. \emph{Journal of the Royal Statistical Society: Series A} 140(4), 411--431.}
\nmref{data}{Data: the lynx counts from PollisDatasets (\fn{lynxmackenzie}); the sunspot numbers as distributed with R (\fn{sunspot.year}), read from R's source and not saved, since SILSO's licence does not allow them to be distributed with Pollis.}
}

\nmcredit
\end{document}
